\begin{lexlean}{Production.ThreatModel}
\useglossary{lexlean.std.bool@1.2.0}
\useglossary{lexlean.std.nat@1.2.0}
\useglossary{prism.arch@1.0.0}
\useglossary{prism.sec@1.0.0}
\importmodule{Foundation.Sec}
\title{Security risk model}
\begin{semanticmodule}
inductive{
  name: Stride,
  type_parameters: [],
  parameters: [],
  constructors: [
    {name: spoofing, fields: []},
    {name: tampering, fields: []},
    {name: repudiation, fields: []},
    {name: informationDisclosure, fields: []},
    {name: denialOfService, fields: []},
    {name: elevationOfPrivilege, fields: []}
  ]
}

inductive{
  name: EvidenceKind,
  type_parameters: [],
  parameters: [],
  constructors: [{name: gate, fields: []}, {name: proof, fields: []}, {name: review, fields: []}]
}

structure{
  name: Evidence,
  type_parameters: [],
  parameters: [],
  fields: [{name: kind, type: EvidenceKind}, {name: reference, type: String}]
}

structure{
  name: ControlSet,
  type_parameters: [],
  parameters: [],
  fields: [
    {name: title, type: String},
    {name: text, type: String},
    {name: version, type: String},
    {name: lastModified, type: String}
  ]
}

structure{
  name: ModeledAsset,
  type_parameters: [],
  parameters: [],
  fields: [{name: asset, type: String}, {name: name, type: String}, {name: text, type: String}]
}

structure{
  name: ModeledThreat,
  type_parameters: [],
  parameters: [],
  fields: [
    {name: threat, type: String},
    {name: name, type: String},
    {name: text, type: String},
    {name: category, type: Stride}
  ]
}

structure{
  name: ModeledControl,
  type_parameters: [],
  parameters: [],
  fields: [
    {name: control, type: String},
    {name: name, type: String},
    {name: text, type: String},
    {name: evidence, type: (List Evidence)}
  ]
}

structure{
  name: ThreatMapping,
  type_parameters: [],
  parameters: [],
  fields: [{name: threat, type: String}, {name: controls, type: (List String)}]
}

structure{
  name: AcceptedRisk,
  type_parameters: [],
  parameters: [],
  fields: [{name: threat, type: String}, {name: rationale, type: String}]
}

structure{
  name: NormalizedThreatModel,
  type_parameters: [],
  parameters: [],
  fields: [
    {name: threatCount, type: Nat},
    {name: categoryCounts, type: (List Nat)},
    {name: mitigationCounts, type: (List Nat)},
    {name: acceptanceCounts, type: (List Nat)},
    {name: evidenceCounts, type: (List Nat)}
  ]
}

definition{
  name: threatListLength,
  parameters: [values : (List Nat)],
  result: Nat,
  body: (match values | List.nil => 0 | List.cons value rest => ((threatListLength rest) + 1)),
  recursive_argument: values
}

definition{
  name: threatAllPositive,
  parameters: [values : (List Nat)],
  result: Bool,
  body: (match values
    | List.nil => true
    | List.cons value rest => ((0 <? value) && (threatAllPositive rest))),
  recursive_argument: values
}

definition{
  name: validateThreatsAnswered,
  parameters: [mitigations : (List Nat), acceptances : (List Nat)],
  result: Bool,
  body: (match mitigations
    | List.nil => true
    | List.cons mitigation rest => (match acceptances
      | List.nil => ((0 <? mitigation) && (validateThreatsAnswered rest acceptances))
      | List.cons acceptance others => (((0 <? mitigation) || (0 <? acceptance)) && (validateThreatsAnswered rest others)))),
  recursive_argument: mitigations
}

definition{
  name: validateThreatCategories,
  parameters: [categoryCounts : (List Nat)],
  result: Bool,
  body: (threatAllPositive categoryCounts)
}

definition{
  name: validateControlEvidence,
  parameters: [evidenceCounts : (List Nat)],
  result: Bool,
  body: (threatAllPositive evidenceCounts)
}

definition{
  name: validateNormalizedThreatModel,
  parameters: [model : NormalizedThreatModel],
  result: Bool,
  body: (((threatListLength (project model categoryCounts)) =? (project model threatCount)) && (((threatListLength (project model mitigationCounts)) =? (project model threatCount)) && (((threatListLength (project model acceptanceCounts)) =? (project model threatCount)) && ((validateThreatCategories (project model categoryCounts)) && ((validateThreatsAnswered (project model mitigationCounts) (project model acceptanceCounts)) && (validateControlEvidence (project model evidenceCounts)))))))
}

definition{
  name: normalizedThreatModelValid,
  parameters: [model : NormalizedThreatModel],
  result: Prop,
  body: ((validateNormalizedThreatModel model) = true)
}

theorem{
  name: normalizedThreatModel_sound_complete,
  parameters: [model : NormalizedThreatModel],
  statement: (((validateNormalizedThreatModel model) = true) <-> (normalizedThreatModelValid model)),
  proof: reflexivity{}
}

definition{
  name: threatModelVector,
  parameters: [
    categoryCounts : (List Nat),
    mitigationCounts : (List Nat),
    acceptanceCounts : (List Nat),
    evidenceCounts : (List Nat)
  ],
  result: NormalizedThreatModel,
  body: (record NormalizedThreatModel {threatCount := 2, categoryCounts := categoryCounts, mitigationCounts := mitigationCounts, acceptanceCounts := acceptanceCounts, evidenceCounts := evidenceCounts})
}

theorem{
  name: threatModelVectorAccepted,
  parameters: [],
  statement: ((validateNormalizedThreatModel (threatModelVector (cons 1 (cons 1 (nil Nat))) (cons 2 (cons 0 (nil Nat))) (cons 0 (cons 1 (nil Nat))) (cons 1 (cons 2 (nil Nat))))) = true),
  proof: decide{}
}

theorem{
  name: unmappedThreatRefused,
  parameters: [],
  statement: ((validateNormalizedThreatModel (threatModelVector (cons 1 (cons 1 (nil Nat))) (cons 2 (cons 0 (nil Nat))) (cons 0 (cons 0 (nil Nat))) (cons 1 (cons 2 (nil Nat))))) = false),
  proof: decide{}
}

theorem{
  name: controlWithoutEvidenceRefused,
  parameters: [],
  statement: ((validateNormalizedThreatModel (threatModelVector (cons 1 (cons 1 (nil Nat))) (cons 2 (cons 0 (nil Nat))) (cons 0 (cons 1 (nil Nat))) (cons 1 (cons 0 (nil Nat))))) = false),
  proof: decide{}
}

theorem{
  name: threatWithoutCategoryRefused,
  parameters: [],
  statement: ((validateNormalizedThreatModel (threatModelVector (cons 1 (cons 0 (nil Nat))) (cons 2 (cons 0 (nil Nat))) (cons 0 (cons 1 (nil Nat))) (cons 1 (cons 2 (nil Nat))))) = false),
  proof: decide{}
}
\end{semanticmodule}
\end{lexlean}
